
\documentclass{article}
\usepackage{colortbl}
\usepackage{makecell}
\usepackage{multirow}
\usepackage{supertabular}

\begin{document}

\newcounter{utterance}

\centering \large Interaction Transcript for game `eqbench', experiment `eq\_bench', episode 131 with reflexion{-}Qwen3.5{-}9B.
\vspace{24pt}

{ \footnotesize  \setcounter{utterance}{1}
\setlength{\tabcolsep}{0pt}
\begin{supertabular}{c@{$\;$}|p{.15\linewidth}@{}p{.15\linewidth}p{.15\linewidth}p{.15\linewidth}p{.15\linewidth}p{.15\linewidth}}
    \# & \multicolumn{2}{c}{Player} && \multicolumn{2}{c}{Game Master} \\
    \hline

    \theutterance \stepcounter{utterance}  
    & & & \multicolumn{4}{p{0.6\linewidth}}{
        \cellcolor[rgb]{0.9,0.9,0.9}{
            \makecell[{{p{\linewidth}}}]{
                \texttt{\tiny{[P1$\langle$GM]}}
                \texttt{You have attempted this same task before and did not succeed. These are the plans you wrote after those attempts.} \\
\texttt{Trial \#1: New plan: I failed because I assigned a high score to "Skeptical" (8) and "Vindicated" (6) while underestimating "Pacified" (2). In the dialogue, Arthur explicitly accepts Sam's anger ("I messed up, I know") and offers a sincere apology without making excuses, which directly addresses Sam's grievance about the neon paint and the wedding speech. This de{-}escalation and validation of Sam's feelings should logically lead to a sense of being heard and soothed, making "Pacified" the dominant emotion. "Skeptical" is less likely because Arthur stops defending his actions and admits fault, and "Vindicated" is incorrect because Sam is not proving a point but rather expressing hurt. Next time, I will prioritize "Pacified" as the highest score (e.g., 9) and lower "Skeptical" to a low single digit (e.g., 2), reflecting that a genuine apology typically resolves the immediate conflict rather than deepening doubt.} \\
\\ 
\\ 
\\ 
\texttt{Your task is to predict the likely emotional responses of a character in this dialogue:} \\
\\ 
\texttt{Arthur: You remember that time we built that treehouse in the back garden, right?} \\
\texttt{Sam: What's that got to do with anything?} \\
\texttt{Arthur: I'm just saying, we were a team, you and me. We built something out of nothing. } \\
\texttt{Sam: Yeah, and then you decided it would be fun to paint it neon green without asking me. } \\
\texttt{Arthur: It was a creative decision, Sammy. } \\
\texttt{Sam: It was a stupid decision, Arthur. Just like your speech at my wedding.} \\
\texttt{Arthur: That's not fair, Sam. I messed up, I know. But I'm trying to apologize here.} \\
\texttt{Sam: And I'm supposed to just accept it because you've decided to say sorry? } \\
\texttt{Arthur: No, not because I've decided to. Because I mean it, Sam. I'm sorry.} \\
\texttt{{[}End dialogue{]}} \\
\\ 
\texttt{At the end of this dialogue, Sam would feel...} \\
\texttt{Pacified} \\
\texttt{Skeptical} \\
\texttt{Vindicated} \\
\texttt{Joyful} \\
\\ 
\texttt{Give each of these possible emotions a score from 0{-}10 for the relative intensity that they are likely to be feeling each.} \\
\\ 
\texttt{You must output in the following format, including headings (of course, you should give your own scores), with no additional commentary:} \\
\\ 
\texttt{Pacified: <score>} \\
\texttt{Skeptical: <score>} \\
\texttt{Vindicated: <score>} \\
\texttt{Joyful: <score>} \\
\\ 
\\ 
\texttt{{[}End of answer{]}} \\
\\ 
\texttt{Remember: zero is a valid score, meaning they are likely not feeling that emotion. You must score at least one emotion > 0.} \\
\\ 
\texttt{Your answer:} \\
            }
        }
    }
     \\ \\

    \theutterance \stepcounter{utterance}  
    & \multicolumn{4}{p{0.6\linewidth}}{
        \cellcolor[rgb]{0.561,0.737,0.561}{
            \makecell[{{p{\linewidth}}}]{
                \texttt{\tiny{[P1$\rangle$GM]}}
                \texttt{Pacified: 2} \\
\texttt{Skeptical: 8} \\
\texttt{Vindicated: 6} \\
\texttt{Joyful: 0} \\
            }
        }
    }
    & & \\ \\

    \theutterance \stepcounter{utterance}  
    & & & \multicolumn{2}{p{0.3\linewidth}}{
        \cellcolor[rgb]{0.9,0.9,0.9}{
            \makecell[{{p{\linewidth}}}]{
                \texttt{\tiny{[GM$|$GM]}}
                \texttt{Pacified: 2} \\
\texttt{Skeptical: 8} \\
\texttt{Vindicated: 6} \\
\texttt{Joyful: 0} \\
            }
        }
    }
    & & \\ \\

    \theutterance \stepcounter{utterance}  
    & & & \multicolumn{2}{p{0.3\linewidth}}{
        \cellcolor[rgb]{0.9,0.9,0.9}{
            \makecell[{{p{\linewidth}}}]{
                \texttt{\tiny{[GM$|$GM]}}
                \texttt{Pacified: 4} \\
\texttt{Skeptical: 7} \\
\texttt{Vindicated: 3} \\
\texttt{Joyful: 0} \\
            }
        }
    }
    & & \\ \\

    \theutterance \stepcounter{utterance}  
    & & & \multicolumn{2}{p{0.3\linewidth}}{
        \cellcolor[rgb]{0.9,0.9,0.9}{
            \makecell[{{p{\linewidth}}}]{
                \texttt{\tiny{[GM$|$GM]}}
                \texttt{game\_result = LOSE} \\
            }
        }
    }
    & & \\ \\

    \theutterance \stepcounter{utterance}  
    & & \multicolumn{4}{p{0.6\linewidth}}{
        \cellcolor[rgb]{0.95,0.95,0.95}{
            \makecell[{{p{\linewidth}}}]{
                \texttt{\tiny{[P1$\leftrightarrow$P]}}
                \texttt{Post{-}episode Reflexion (not sent to GM):} \\
\texttt{New plan: I failed because I misinterpreted the emotional trajectory of the dialogue. Arthur's apology, while sincere, was immediately followed by Sam's rejection ("I'm supposed to just accept it?"), which indicates the conflict is unresolved and Sam's anger is still active. Therefore, "Pacified" is incorrect as the primary outcome. Sam feels "Vindicated" because Arthur's admission of fault regarding the neon paint and the wedding speech validates her previous criticism, confirming her stance that his actions were wrong. "Skeptical" is also present but secondary to the strong validation of her grievances. Next time, I will assign the highest score to "Vindicated" (e.g., 9) to reflect that Sam's anger is fueled by Arthur's admission of past wrongs, and give a moderate score to "Skeptical" (e.g., 4) to acknowledge the tension, while keeping "Pacified" and "Joyful" at 0.} \\
            }
        }
    }
    & \\ \\

\end{supertabular}
}

\end{document}
